\begin{abstract}
This thesis studies how exposure to international maritime trade shapes subnational conflict in Nigeria. For doing so, we built an exposure to trade index. We interacted changes in container port throughput (CPT) with the Liner Shipping Connectivity Index (LSCI) and the inverse distance to the nearest commercial port. Then, we studies the effect of the trade index to conflict, in particular ADM2 level events from ACLED (2009–2019). Using two-way fixed effects we found that higher exposure to trade is associated with a stabilizing effect on conflicts at the micro level, especially in poorer areas, while effects are near zero or weaker in more developed districts. Accounting for heterogeneity, results reported a stronger effect in the eastern coastal regions and for ports oriented toward vegetable goods and imports, and they attenuate in regions with abundance of water availability, a proxy for agricultural rents, aligning with the result of the baseline model. We continued the analysis estimating a falsification test using non-commercial ports and that supported our main hypothesis, demonstrating an effect of trade in conflict, rather than mere geographic proximity. The findings highlights the central role of pre-existing economic and spatial inequalities in mediating the stabilizing effects of globalization

\end{abstract}